% Options for packages loaded elsewhere
\PassOptionsToPackage{unicode}{hyperref}
\PassOptionsToPackage{hyphens}{url}
\documentclass[
]{article}
\usepackage{xcolor}
\usepackage{amsmath,amssymb}
\setcounter{secnumdepth}{-\maxdimen} % remove section numbering
\usepackage{iftex}
\ifPDFTeX
  \usepackage[T1]{fontenc}
  \usepackage[utf8]{inputenc}
  \usepackage{textcomp} % provide euro and other symbols
\else % if luatex or xetex
  \usepackage{unicode-math} % this also loads fontspec
  \defaultfontfeatures{Scale=MatchLowercase}
  \defaultfontfeatures[\rmfamily]{Ligatures=TeX,Scale=1}
\fi
\usepackage{lmodern}
\ifPDFTeX\else
  % xetex/luatex font selection
\fi
% Use upquote if available, for straight quotes in verbatim environments
\IfFileExists{upquote.sty}{\usepackage{upquote}}{}
\IfFileExists{microtype.sty}{% use microtype if available
  \usepackage[]{microtype}
  \UseMicrotypeSet[protrusion]{basicmath} % disable protrusion for tt fonts
}{}
\makeatletter
\@ifundefined{KOMAClassName}{% if non-KOMA class
  \IfFileExists{parskip.sty}{%
    \usepackage{parskip}
  }{% else
    \setlength{\parindent}{0pt}
    \setlength{\parskip}{6pt plus 2pt minus 1pt}}
}{% if KOMA class
  \KOMAoptions{parskip=half}}
\makeatother
\usepackage{graphicx}
\makeatletter
\newsavebox\pandoc@box
\newcommand*\pandocbounded[1]{% scales image to fit in text height/width
  \sbox\pandoc@box{#1}%
  \Gscale@div\@tempa{\textheight}{\dimexpr\ht\pandoc@box+\dp\pandoc@box\relax}%
  \Gscale@div\@tempb{\linewidth}{\wd\pandoc@box}%
  \ifdim\@tempb\p@<\@tempa\p@\let\@tempa\@tempb\fi% select the smaller of both
  \ifdim\@tempa\p@<\p@\scalebox{\@tempa}{\usebox\pandoc@box}%
  \else\usebox{\pandoc@box}%
  \fi%
}
% Set default figure placement to htbp
\def\fps@figure{htbp}
\makeatother
\setlength{\emergencystretch}{3em} % prevent overfull lines
\providecommand{\tightlist}{%
  \setlength{\itemsep}{0pt}\setlength{\parskip}{0pt}}
\usepackage{bookmark}
\IfFileExists{xurl.sty}{\usepackage{xurl}}{} % add URL line breaks if available
\urlstyle{same}
% fallback for those not using the hyperref driver hyperxmp:
\makeatletter
\@ifundefined{xmpquote}{\newcommand{\xmpquote}[1]{#1}}{}
\makeatother
\hypersetup{
  hidelinks,
  pdfcreator={LaTeX via pandoc}}

\author{}
\date{}

\begin{document}

\section{Artefact 1 --- Clean Database \& Univariate Description
Report}\label{artefact-1-clean-database-univariate-description-report}

\begin{quote}
\textbf{Course}: Quantitative Data Analysis (M1 -- S7) \textbf{Course
Code}: \texttt{2627\_ECO\_2\_EN\_009} --- Academic Year 2026--2027
\textbf{Institution}: EM Normandie Business School \textbf{Lecturer}:
Dr.~NGUYEN Anh-Tuan \textbf{Deliverable}: \textbf{Artefact 1 (20 Marks)}
\textbf{Accompanying Files}: \texttt{sp500\_esg\_jamovi.csv},
\texttt{sp500\_esg\_jamovi.xlsx}, \texttt{Approval\_Slip.md},
\texttt{Data\_Cleaning\_Log.md}
\end{quote}

\begin{center}\rule{0.5\linewidth}{0.5pt}\end{center}

\subsection{1. Submission Overview \& Dataset
Access}\label{submission-overview-dataset-access}

This report constitutes the official submission for \textbf{Artefact 1},
fulfilling all requirements specified in Slide 41 of the course
syllabus. The underlying dataset comprises \(N = 503\) publicly traded
corporations from the S\&P 500 index for Fiscal Year 2025 (FY2025).

\subsubsection{Available Data Files}\label{available-data-files}

\begin{enumerate}
\def\labelenumi{\arabic{enumi}.}
\tightlist
\item
  \textbf{Jamovi Ready CSV}:
  \href{file:///Users/david/Library/Mobile\%20Documents/com~apple~CloudDocs/Etudes\%20Supp/EMN\%20M1\%20/M1\%20-\%20S7/Data\%20Analysis\%20Quantitative/Group\%20Work\%20S&P\%20500/data/processed/sp500_esg_jamovi.csv}{\texttt{data/processed/sp500\_esg\_jamovi.csv}}

  \begin{itemize}
  \tightlist
  \item
    Formatted with standardized short headers, no spaces, no accents,
    and period decimals. Ready for instant import via Jamovi
    (\texttt{≡\ menu\ →\ Open\ →\ Browse}).
  \end{itemize}
\item
  \textbf{Jamovi / Excel Master Workbook}:
  \href{file:///Users/david/Library/Mobile\%20Documents/com~apple~CloudDocs/Etudes\%20Supp/EMN\%20M1\%20/M1\%20-\%20S7/Data\%20Analysis\%20Quantitative/Group\%20Work\%20S&P\%20500/data/processed/sp500_esg_jamovi.xlsx}{\texttt{data/processed/sp500\_esg\_jamovi.xlsx}}

  \begin{itemize}
  \tightlist
  \item
    Includes two dedicated sheets: \texttt{Data} (the full analytical
    matrix) and \texttt{Codebook\_Jamovi} (detailed metadata with
    variable definitions and measurement types).
  \end{itemize}
\end{enumerate}

\begin{center}\rule{0.5\linewidth}{0.5pt}\end{center}

\subsection{2. Summary of Provenance \& Approval Slip (Slide
35)}\label{summary-of-provenance-approval-slip-slide-35}

\begin{itemize}
\tightlist
\item
  \textbf{Dataset Title}: S\&P 500 Corporate Governance \& Financial
  Performance Dataset (FY2025).
\item
  \textbf{Producers}: Standard \& Poor's, Institutional Shareholder
  Services (ISS), and Yahoo Finance (SEC 10-K audited filings).
\item
  \textbf{Unit of Observation}: 1 row = 1 publicly traded S\&P 500 firm
  (\(N = 503\)).
\item
  \textbf{Research Question}: \emph{``Does corporate governance risk
  (ISS QualityScore) compromise financial profitability and market
  valuation among S\&P 500 firms, or does strong governance provide a
  resilience and performance premium?''}
\item
  \textbf{Complete Details}: Refer to
  \href{file:///Users/david/Library/Mobile\%20Documents/com~apple~CloudDocs/Etudes\%20Supp/EMN\%20M1\%20/M1\%20-\%20S7/Data\%20Analysis\%20Quantitative/Group\%20Work\%20S&P\%20500/docs/Approval_Slip.md}{\texttt{docs/Approval\_Slip.md}}.
\end{itemize}

\begin{center}\rule{0.5\linewidth}{0.5pt}\end{center}

\subsection{3. Data Cleaning Summary (Slide
39)}\label{data-cleaning-summary-slide-39}

Prior to statistical computation, the raw extract underwent systematic
verification against the 6 student traps (Slide 39):

\begin{enumerate}
\def\labelenumi{\arabic{enumi}.}
\tightlist
\item
  \textbf{Fiserv Remediation}: Ticker updated from \texttt{FISV} to
  \texttt{FI} to resolve a missing observation caused by its 2023 NYSE
  transfer, restoring full financial and governance coverage.
\item
  \textbf{Zero Duplicates}: Verified 100\% uniqueness across all 503
  equity tickers.
\item
  \textbf{Negative Equity Treatment}: Identified 32 firms with undefined
  \texttt{ROE} due to multi-year share buybacks driving book equity
  below zero (e.g.~AbbVie, Altria, Domino's). Documented transparently
  without dropping observations, using \texttt{Profit\_Margin} as the
  universal profitability metric.
\item
  \textbf{Outlier Documentation}: Retained Masco Corp's authentic
  \(58.625\) ROE while establishing the Median (\(16.8\%\)) as the true
  managerial benchmark to avoid mean distortion (Slide 21).
\item
  \textbf{Sector Consolidation}: Grouped 11 GICS sectors into exactly 5
  nominal categories (\(n = 60\) to \(148\)).
\item
  \textbf{Full Audit}: Refer to
  \href{file:///Users/david/Library/Mobile\%20Documents/com~apple~CloudDocs/Etudes\%20Supp/EMN\%20M1\%20/M1\%20-\%20S7/Data\%20Analysis\%20Quantitative/Group\%20Work\%20S&P\%20500/docs/Data_Cleaning_Log.md}{\texttt{docs/Data\_Cleaning\_Log.md}}.
\end{enumerate}

\begin{center}\rule{0.5\linewidth}{0.5pt}\end{center}

\subsection{4. Univariate Description of Five Variables (Slide
41)}\label{univariate-description-of-five-variables-slide-41}

In accordance with course directives (\emph{``Half the marks are for
points 4--6 --- the interpretation. Producing output is not analysis.
Explaining it is''}), the following five variables are profiled in
complete sentences adhering to statistical protocol.

\subsubsection{\texorpdfstring{4.1 Variable 1: \texttt{Sector} (Nominal
Variable --- 5
Categories)}{4.1 Variable 1: Sector (Nominal Variable --- 5 Categories)}}\label{variable-1-sector-nominal-variable-5-categories}

\begin{itemize}
\tightlist
\item
  \textbf{Nature}: Qualitative nominal factor representing the primary
  industry classification of the firm. Because categories possess no
  mathematical rank or numerical magnitude, computing a mean or median
  is strictly invalid (Slide 19); the distribution must be described
  exclusively through frequency tables and the \textbf{mode} (Slide 20).
\item
  \textbf{Distribution \& Mode}: Across the \(N = 503\) firms, the modal
  category is \textbf{Industrials \& Energy}, representing \(148\) firms
  (\(29.42\%\)), followed by \textbf{Consumer} with \(116\) firms
  (\(23.06\%\)), \textbf{Tech \& Comms} with \(108\) firms
  (\(21.47\%\)), \textbf{Finance} with \(71\) firms (\(14.12\%\)), and
  \textbf{Healthcare} with \(60\) firms (\(11.93\%\)).
\item
  \textbf{Managerial Interpretation}: The S\&P 500 is structurally
  anchored in industrial, cyclical, and technological infrastructure
  rather than consumer staples or pure healthcare. Every sector retains
  at least 60 observations, satisfying statistical power requirements
  for the upcoming One-Way ANOVA in Session 3.
\end{itemize}

\begin{verbatim}
Sector Frequency Breakdown:
1. Industrials & Energy : 148 (29.4%) [MODE]
2. Consumer             : 116 (23.1%)
3. Tech & Comms         : 108 (21.5%)
4. Finance              :  71 (14.1%)
5. Healthcare           :  60 (11.9%)
Total Valid Observations: 503 (100.0%)
\end{verbatim}

\begin{figure}
\centering
\pandocbounded{\includegraphics[keepaspectratio,alt={Figure 1: S\&P 500 Sector Distribution across 5 Consolidated Categories}]{figures/sector_distribution.png}}
\caption{Figure 1: S\&P 500 Sector Distribution across 5 Consolidated
Categories}
\end{figure}

\begin{center}\rule{0.5\linewidth}{0.5pt}\end{center}

\subsubsection{\texorpdfstring{4.2 Variable 2:
\texttt{Governance\_Risk\_Level} (Ordinal Variable --- 3
Tiers)}{4.2 Variable 2: Governance\_Risk\_Level (Ordinal Variable --- 3 Tiers)}}\label{variable-2-governance_risk_level-ordinal-variable-3-tiers}

\begin{itemize}
\tightlist
\item
  \textbf{Nature}: Qualitative ordinal factor classifying firms into
  three ordered governance risk tiers derived from the ISS Overall
  QualityScore: \emph{Low} (scores 1--3, best governance), \emph{Medium}
  (scores 4--6), and \emph{High} (scores 7--10, poorest governance).
  While categories follow a natural sequence, intervals between tiers
  are not mathematically identical (Slide 13); central tendency is
  captured by the mode and median rank.
\item
  \textbf{Distribution \& Central Tendency}: Among the \(496\) reporting
  corporations (\(7\) firms unrated by ISS), the modal tier is
  \textbf{High Governance Risk}, encompassing \(197\) firms
  (\(39.72\%\)), while \textbf{Low Governance Risk} accounts for \(150\)
  firms (\(30.24\%\)) and \textbf{Medium Governance Risk} accounts for
  \(149\) firms (\(30.04\%\)). The median firm falls exactly in the
  \textbf{Medium} tier.
\item
  \textbf{Managerial Interpretation}: Contrary to common assumptions,
  nearly \(40\%\) of S\&P 500 corporations exhibit high institutional
  governance risk scores, predominantly driven by contentious executive
  compensation packages and shareholder rights restrictions. This
  substantial variance provides an ideal foundation for testing risk
  differentials in Session 2 (Chi-Square) and Session 4 (Regression).
\end{itemize}

\begin{figure}
\centering
\pandocbounded{\includegraphics[keepaspectratio,alt={Figure 2: Distribution of Institutional Governance Risk Tiers}]{figures/governance_risk_distribution.png}}
\caption{Figure 2: Distribution of Institutional Governance Risk Tiers}
\end{figure}

\begin{center}\rule{0.5\linewidth}{0.5pt}\end{center}

\subsubsection{\texorpdfstring{4.3 Variable 3: \texttt{Profit\_Margin}
(Continuous Variable --- Ratio
Scale)}{4.3 Variable 3: Profit\_Margin (Continuous Variable --- Ratio Scale)}}\label{variable-3-profit_margin-continuous-variable-ratio-scale}

\begin{itemize}
\tightlist
\item
  \textbf{Nature}: Quantitative continuous ratio variable measuring net
  income as a proportion of total revenue for FY2025. Because zero
  represents a true economic threshold (break-even), ratios and
  parametric moments are mathematically defined (Slide 14).
\item
  \textbf{Statistical Profile (\(N = 503\))}:

  \begin{itemize}
  \tightlist
  \item
    \emph{Central Tendency}: Mean \(\bar{x} = 0.1479\) (\(14.79\%\)),
    Median \(\text{Mdn} = 0.1315\) (\(13.15\%\)), Rounded Mode
    \(= 0.1100\) (\(11.00\%\)).
  \item
    \emph{Dispersion}: Standard Deviation \(s = 0.1856\) (\(18.56\%\)),
    Variance \(s^2 = 0.0345\), Interquartile Range \(IQR = 0.1446\)
    (\(14.46\%\), spanning \(Q_1 = 7.50\%\) to \(Q_3 = 21.97\%\)). Full
    range spans from \(-2.3010\) (Lumentum Holdings) to \(+0.7295\)
    (Western Digital).
  \item
    \emph{Shape}: Fisher-Pearson Skewness \(g_1 = -5.4205\) (severe
    negative/left skewness), Excess Kurtosis \(g_2 = 69.9129\) (extreme
    leptokurtosis with heavy loss-making tails).
  \end{itemize}
\item
  \textbf{Managerial Interpretation}: The median S\&P 500 company
  generates \(13.15\) cents of bottom-line profit for every dollar of
  turnover. However, the substantial standard deviation (\(18.56\%\))
  relative to the mean confirms pronounced profit dispersion across
  business models, with a small cluster of tech/semiconductor companies
  posting massive margins while distressed or restructuring firms pull
  down the arithmetic average.
\end{itemize}

\begin{center}\rule{0.5\linewidth}{0.5pt}\end{center}

\subsubsection{\texorpdfstring{4.4 Variable 4: \texttt{ROE} (Continuous
Variable --- Ratio
Scale)}{4.4 Variable 4: ROE (Continuous Variable --- Ratio Scale)}}\label{variable-4-roe-continuous-variable-ratio-scale}

\begin{itemize}
\tightlist
\item
  \textbf{Nature}: Quantitative continuous variable measuring annual net
  earnings generated per dollar of book equity.
\item
  \textbf{Statistical Profile (\(N = 471\) valid, \(32\) negative-equity
  firms)}:

  \begin{itemize}
  \tightlist
  \item
    \emph{Central Tendency}: Mean \(\bar{x} = 0.3952\) (\(39.52\%\)),
    Median \(\text{Mdn} = 0.1676\) (\(16.76\%\)), Rounded Mode
    \(= 0.1200\) (\(12.00\%\)).
  \item
    \emph{Dispersion}: Standard Deviation \(s = 2.7231\) (\(272.31\%\)),
    Interquartile Range \(IQR = 0.2195\) (\(21.95\%\), \(Q_1 = 9.32\%\)
    to \(Q_3 = 31.27\%\)). Range spans from \(-2.4003\) to \(+58.6250\)
    (Masco Corp).
  \item
    \emph{Shape}: Skewness \(g_1 = 20.9053\), Kurtosis
    \(g_2 = 447.6009\).
  \end{itemize}
\item
  \textbf{Managerial Interpretation}: The massive divergence between the
  arithmetic mean (\(39.52\%\)) and the median (\(16.76\%\)) illustrates
  the textbook danger highlighted in Slide 21 (\emph{``The mean is
  destroyed by outliers. The median is robust''}). Masco Corp's
  near-zero equity balance inflates its ROE to \(5,862.5\%\), heavily
  distorting the mean. A manager analyzing corporate performance must
  rely on the median ROE of \(16.76\%\) as the true representative
  baseline.
\end{itemize}

\begin{center}\rule{0.5\linewidth}{0.5pt}\end{center}

\subsubsection{\texorpdfstring{4.5 Variable 5: \texttt{Market\_Cap\_B}
(Continuous Variable --- Ratio
Scale)}{4.5 Variable 5: Market\_Cap\_B (Continuous Variable --- Ratio Scale)}}\label{variable-5-market_cap_b-continuous-variable-ratio-scale}

\begin{itemize}
\tightlist
\item
  \textbf{Nature}: Quantitative continuous ratio variable measuring the
  total equity market valuation in billions of USD.
\item
  \textbf{Statistical Profile (\(N = 503\))}:

  \begin{itemize}
  \tightlist
  \item
    \emph{Central Tendency}: Mean \(\bar{x} = \$147.60\text{B}\), Median
    \(\text{Mdn} = \$44.47\text{B}\), Mode Cluster
    \(\approx \$20.00\text{B}\).
  \item
    \emph{Dispersion}: Standard Deviation \(s = \$490.01\text{B}\),
    Interquartile Range \(IQR = \$71.79\text{B}\)
    (\(Q_1 = \$22.21\text{B}\) to \(Q_3 = \$94.00\text{B}\)). Range
    spans from \(\$6.48\text{B}\) to \(\$5,367.15\text{B}\).
  \item
    \emph{Shape}: Skewness \(g_1 = 7.8532\), Kurtosis \(g_2 = 68.1146\).
  \end{itemize}
\item
  \textbf{Managerial Interpretation}: S\&P 500 capitalization exhibits
  acute right-skewness (\emph{Pareto distribution}): while half the
  index is valued below \(\$44.5\text{B}\), a handful of multi-trillion
  mega-cap tech conglomerates (NVIDIA, Apple, Microsoft, Alphabet,
  Amazon) inflate the average valuation to \(\$147.6\text{B}\). Size
  must therefore be modeled via quartiles
  (\texttt{Market\_Cap\_Quartile}) or logarithmic transformations in
  subsequent regression modeling.
\end{itemize}

\begin{figure}
\centering
\pandocbounded{\includegraphics[keepaspectratio,alt={Figure 3: S\&P 500 Market Capitalization Distribution and Pareto Skewness}]{figures/market_cap_distribution.png}}
\caption{Figure 3: S\&P 500 Market Capitalization Distribution and
Pareto Skewness}
\end{figure}

\begin{center}\rule{0.5\linewidth}{0.5pt}\end{center}

\subsection{\texorpdfstring{5. Normality Evaluation of a Continuous
Variable
(\texttt{Profit\_Margin})}{5. Normality Evaluation of a Continuous Variable (Profit\_Margin)}}\label{normality-evaluation-of-a-continuous-variable-profit_margin}

In strict accordance with Slide 24 (\emph{``Three converging pieces of
evidence --- use all three, never just one''}), we evaluate whether
\texttt{Profit\_Margin} satisfies the assumption of normality:

\begin{verbatim}
========================================================================================
NORMALITY DIAGNOSTIC MATRIX: Net Profit Margin (N = 503)
========================================================================================
Test Component        Empirical Value            Normality Criterion           Verdict
----------------------------------------------------------------------------------------
1. Three Centres      Mean   = 0.1479            Mean ≈ Median ≈ Mode          VIOLATED
                      Median = 0.1315            (Centres diverge by > 0.035)
                      Mode   = 0.1100

2. Shape Parameters   Skewness = -5.42           Both strictly in [-1, +1]     VIOLATED
                      Kurtosis = 69.91           (Severe left skew & leptokurtic)

3. Formal Test        Shapiro-Wilk W = 0.6530    p > 0.05 (Fail to reject H0)  REJECTED
                      p-value < .001                                            (p < .001)
========================================================================================
FINAL VERDICT: Severe violation of normality. Parametric normality assumption REJECTED.
========================================================================================
\end{verbatim}

\begin{figure}
\centering
\pandocbounded{\includegraphics[keepaspectratio,alt={Figure 4: Net Profit Margin Normality Diagnostics --- Histogram, Density, and Q-Q Plot}]{figures/profit_margin_normality.png}}
\caption{Figure 4: Net Profit Margin Normality Diagnostics ---
Histogram, Density, and Q-Q Plot}
\end{figure}

\subsubsection{The Three Converging Pieces of Evidence
Explained:}\label{the-three-converging-pieces-of-evidence-explained}

\begin{enumerate}
\def\labelenumi{\arabic{enumi}.}
\tightlist
\item
  \textbf{Evidence 1: Comparison of the Three Centres (Slide 24)} In a
  normal distribution, the curve is perfectly symmetric and
  \(\text{Mean} \approx \text{Median} \approx \text{Mode}\). Here, the
  arithmetic mean (\(0.1479\)) noticeably exceeds the median
  (\(0.1315\)), which in turn exceeds the modal cluster (\(0.1100\)).
  This clear divergence indicates positive asymmetry across typical
  operating firms, countered by severe negative tail outliers.
\item
  \textbf{Evidence 2: Skewness and Kurtosis Rule of Thumb (Slide 23 \&
  24)} The course rule of thumb dictates that both skewness and kurtosis
  must fall within \([-1.0, +1.0]\) for a distribution to be considered
  normal. \texttt{Profit\_Margin} yields a skewness of \(-5.42\) and an
  excess kurtosis of \(+69.91\), violently breaching the normative
  boundaries due to extreme loss-making firms in cyclical sectors.
\item
  \textbf{Evidence 3: Graphical Inspection \& Formal Shapiro-Wilk Test
  (Slide 24)}

  \begin{itemize}
  \tightlist
  \item
    \emph{Graphical Inspection}: The density plot demonstrates a sharp,
    narrow peak with long heavy tails, while the Normal Q-Q plot reveals
    severe departures from the 45-degree theoretical reference line at
    both extremities.
  \item
    \emph{Formal Test}: The Shapiro-Wilk test tests the null hypothesis
    \(H_0\): \emph{``The population is normally distributed''}. The test
    statistic is \(W = 0.6530\) with a significance level of
    \(p < .001\). Because \(p < 0.05\), we decisively \textbf{reject
    \(H_0\)}.
  \end{itemize}
\end{enumerate}

\textbf{Methodological Implication for Sessions 2, 3, and 4}: Because
normality is violated, Session 3 group comparisons must employ
\textbf{Welch's ANOVA} or \textbf{Kruskal-Wallis} tests, with post-hoc
comparisons using \textbf{Games-Howell} rather than standard Tukey HSD
whenever variance homogeneity is also breached.

\begin{center}\rule{0.5\linewidth}{0.5pt}\end{center}

\subsection{6. Two-Sentence Managerial Takeaways (Slide
41)}\label{two-sentence-managerial-takeaways-slide-41}

\begin{quote}
\textbf{``Because profit margins and market valuations deviate sharply
from normal distributions with high positive and negative tail skewness,
managers and corporate analysts must evaluate operating performance
against sector medians (\(13.15\%\)) rather than distorted index
averages (\(14.79\%\)). On Monday morning, leadership should benchmark
corporate governance policies not to arbitrary peer averages, but
against the \(39.7\%\) high-risk corporate governance cohort to preempt
investor pushback on executive compensation and board independence.''}
\end{quote}

\begin{center}\rule{0.5\linewidth}{0.5pt}\end{center}

\subsection{7. Jamovi Navigation Guide for Session
1}\label{jamovi-navigation-guide-for-session-1}

To replicate these findings directly in the Jamovi desktop environment
(Slides 36--40):

\begin{enumerate}
\def\labelenumi{\arabic{enumi}.}
\tightlist
\item
  \textbf{Import Data}: Launch Jamovi \(\rightarrow\) Click
  \texttt{≡\ menu} \(\rightarrow\) \texttt{Open} \(\rightarrow\)
  \texttt{Browse} \(\rightarrow\) Select
  \texttt{sp500\_esg\_jamovi.csv}.
\item
  \textbf{Verify Measure Types (Slide 38)}:

  \begin{itemize}
  \tightlist
  \item
    \texttt{Ticker}: Set to \texttt{ID}.
  \item
    \texttt{Sector}: Set to \texttt{Nominal}.
  \item
    \texttt{Market\_Cap\_Quartile} \& \texttt{Governance\_Risk\_Level}:
    Set to \texttt{Ordinal}.
  \item
    \texttt{Total\_Revenue\_B}, \texttt{Market\_Cap\_B},
    \texttt{Profit\_Margin}, \texttt{ROE}, \texttt{Beta},
    \texttt{Overall\_Governance\_Risk}: Set to \texttt{Continuous}.
  \end{itemize}
\item
  \textbf{Run Descriptive Statistics (Slide 40)}:

  \begin{itemize}
  \tightlist
  \item
    Navigate to \texttt{Analyses} \(\rightarrow\) \texttt{Exploration}
    \(\rightarrow\) \texttt{Descriptives}.
  \item
    For continuous variables: Check \texttt{Mean}, \texttt{Median},
    \texttt{Mode}, \texttt{Std.\ deviation}, \texttt{Min}, \texttt{Max},
    \texttt{IQR}, \texttt{Skewness}, \texttt{Kurtosis},
    \texttt{Shapiro-Wilk}, and select \texttt{Histogram},
    \texttt{Density}, and \texttt{Box\ plot}.
  \item
    For categorical variables: Move \texttt{Sector} and
    \texttt{Governance\_Risk\_Level} to Variables \(\rightarrow\) Check
    \texttt{Frequency\ tables} and select \texttt{Bar\ plot}.
  \end{itemize}
\item
  \textbf{Save Session}: Click \texttt{≡\ menu} \(\rightarrow\)
  \texttt{Save\ As} \(\rightarrow\) Save as
  \texttt{sp500\_esg\_session1.omv}.
\end{enumerate}

\end{document}
